\section{Projet 14 — Memory}

\subsection*{Objectif du jeu}
Retrouver les 8 paires de cartes identiques cachées dans une grille de 16 cartes, en un minimum de tentatives.

\subsection*{Règles du jeu}
\begin{itemize}
  \item La grille compte \textbf{16 cartes} ($4\times4$) numérotées de \texttt{1} à \texttt{16}. Chaque carte porte une \textbf{lettre} sur sa face : les lettres \texttt{'A'} à \texttt{'H'}, chacune en \textbf{double} (8 paires).
  \item Les lettres sont attribuées aux cartes par un \textbf{tirage aléatoire sans remise} (comme les cartes du Poker) : on ne les mélange pas, on les tire une à une.
  \item Par défaut toutes les cartes sont \textbf{cachées} : le dos d'une carte affiche un \texttt{?} et, juste en dessous, son \textbf{numéro}, toujours visible.
  \item À chaque tentative, le joueur indique le \textbf{numéro de deux cartes distinctes} ; leurs faces (les lettres) sont révélées.
  \item Si les deux lettres sont \textbf{identiques}, la paire est trouvée et reste visible définitivement ; sinon les deux cartes sont de nouveau cachées après une courte pause (\texttt{sleep()}, voir section Bonus).
  \item Le jeu est \textbf{gagné} lorsque les 8 paires sont trouvées.
\end{itemize}

\subsection*{Affichage de la grille}
Au départ, toutes les cartes sont cachées. L'affichage ci-dessous est fourni dans le squelette :
\begin{lstlisting}[basicstyle=\ttfamily\small,numbers=none]
+----+----+----+----+
| ?  | ?  | ?  | ?  |
|  1 |  2 |  3 |  4 |
+----+----+----+----+
| ?  | ?  | ?  | ?  |
|  5 |  6 |  7 |  8 |
+----+----+----+----+
| ?  | ?  | ?  | ?  |
|  9 | 10 | 11 | 12 |
+----+----+----+----+
| ?  | ?  | ?  | ?  |
| 13 | 14 | 15 | 16 |
+----+----+----+----+
\end{lstlisting}
Une carte \textbf{trouvée} affiche sa lettre à la place du \texttt{?} (et de son numéro).

\textit{Conseil :} on utilise trois tableaux de 16 cases : les \textbf{lettres} des cartes, l'\textbf{état} de chaque carte (trouvée ou non) et un \textbf{paquet} de 16 lettres (\texttt{AABBCC\dots HH}) dans lequel on tire. Le tirage sans remise remplace la case tirée par la dernière case du paquet, puis réduit le paquet d'une case (voir la fiche d'aide « Techniques utiles », section Tirage sans remise). Pour tester, remplacez le tirage aléatoire par un tirage non aléatoire.

\subsection*{Cahier des charges}
\begin{UPSTICahierDesCharges}{1}
\begin{itemize}
  \item Remplir le paquet avec les lettres \texttt{'A'} à \texttt{'H'} \textbf{en double}.
  \item Attribuer une lettre à chacune des 16 cartes par \textbf{tirage aléatoire sans remise} dans le paquet.
  \item Afficher la grille \textbf{cachée} (fournie dans le squelette) ; en mode test, afficher aussi les lettres réelles.
\end{itemize}
\end{UPSTICahierDesCharges}

\begin{UPSTICahierDesCharges}{2}
\begin{itemize}
  \item Saisir les \textbf{numéros de deux cartes} (vérifier qu'ils sont compris entre 1 et 16, distincts et que les cartes ne sont pas déjà trouvées).
  \item Afficher \textbf{temporairement} les deux lettres réelles, puis les \textbf{comparer}.
  \item Mettre à jour le \textbf{tableau d'état} (carte trouvée oui/non) et ré-afficher la grille après la pause.
\end{itemize}
\end{UPSTICahierDesCharges}

\begin{UPSTICahierDesCharges}{3}
\begin{itemize}
  \item Compter et afficher le \textbf{nombre de tentatives}.
  \item Détecter la \textbf{victoire} (8 paires sur 8).
  \item Structurer en fonctions (remplissage du paquet, tirage, affichage, saisie, comparaison).
\end{itemize}
\end{UPSTICahierDesCharges}

\textit{Amélioration :} mode \textbf{2 joueurs} : le joueur qui trouve une paire rejoue ; le vainqueur est celui qui a trouvé le plus de paires.

